%% kapa: exam=2 round=33 year=2022 period=3 subject=compensation_law subject_ko=감정평가 및 보상법규 seq=1 number=1 points=40
%% source: Q-Net 감정평가사 기출문제 (cst00309) artlSeq=5210838 file=raw/1008/5210838/2022년 제33회 감정평가사 2차시험 3교시 문제지(감정평가 및 보상법규).pdf
%% conversion: kordoc 4.24.51 -> tools/gen2.py (auto); review=proofread-gpt-astra; flags=none
\begin{longquestion}{【문제 1】}
X는 「도시 및 주거환경정비법」(이하 ‘도시정비법’이라 함)에 따른 재개발 정비사업조합이고, 甲은 X의 조합원으로서, 해당 정비사업구역 내에 있는 A토지와 B토지의 소유자이다. A토지와 B토지는 연접하고 있고 그 지목이 모두 대(垈)에 해당하지만, A토지는 「사도법」에 따른 사도가 아닌데도 불특정 다수인의 통행에 장기간 제공되어 왔고, B토지는 甲이 소유한 건축물의 부지로서 그 건축물의 일부에 임차인 乙이 거주하고 있다. X는 도시정비법 제72조 제1항에 따라 분양신청기간을 공고하였으나 甲은 그 기간 내에 분양신청을 하지 않았다. 이에 따라 X는 甲을 분양대상자에서 제외하고 관리처분계획을 수립하여 인가를 받았고, 그에 불복하는 행정심판이나 행정소송은 없었다. X는 도시정비법 제73조 제1항에 따른 甲과의 보상협의가 이루어지지 않자 A토지와 B토지에 관하여 관할 토지수용위원회에 수용재결을 신청하였고, 관할 토지수용위원회는 A토지와 B토지를 수용한다는 내용의 수용재결을 하였다. 다음 물음에 답하시오. (40점) \par
\medskip\noindent 물음 1) 甲이 수용재결에 대한 취소소송을 제기하면서, ‘X가 도시정비법 제72조 제1항에 따라 분양신청기간과 그 기간 내에 분양신청을 할 수 있다는 취지를 명백히 표시하여 통지하여야 하는데도 이러한 절차를 제대로 거치지 않았다’고 주장할 경우에, 甲의 주장이 사실이라면 법원은 그것을 이유로 수용재결을 취소할 수 있는지 설명하시오. (단, 사실심 변론종결 전에 도시정비법에 따른 이전고시가 효력을 발생한 경우와 그렇지 않은 경우를 구분하여 설명할 것) (10점) \par
\medskip\noindent 물음 2) 「공익사업을 위한 토지 등의 취득 및 보상에 관한 법률 시행규칙」(이하 ‘토지보상법 시행규칙’이라 함) 제26조 제1항에 따른 ‘사실상의 사도’의 요건을 설명하고, 이에 따라 A토지가 사실상의 사도로 인정되는 경우와 그렇지 않은 경우에 보상기준이 어떻게 달라지는지 설명하시오. (10점) \par
\medskip\noindent 물음 3) 주거이전비에 관하여 甲은 토지보상법 시행규칙 제54조 제1항에 따른 요건을 갖추고 있고, 乙은 같은 조 제2항에 따른 요건을 갖추고 있다. 관할 토지수용위원회는 수용재결을 하면서 甲의 주거이전비에 관하여는 재결을 하였으나 乙의 주거이전비에 관하여는 재결을 하지 않았다. 甲은 주거이전비의 증액을 청구하고자 하고, 乙은 주거이전비의 지급을 청구하고자 한다. 甲과 乙의 권리구제에 적합한 소송을 설명하시오. (20점) \par
\begin{framedtext}
참조 조문 \par
<도시 및 주거환경정비법> \par
제72조(분양공고 및 분양신청) ① 사업시행자는 제50조제9항에 따른 사업시행계획인가의 고시가 있은 날(사업시행계획인가 이후 시공자를 선정한 경우에는 시공자와 계약을 체결한 날)부터 120일 이내에 다음 각 호의 사항을 토지등소유자에게 통지하고, 분양의 대상이 되는 대지 또는 건축물의 내역 등 대통령령으로 정하는 사항을 해당 지역에서 발간되는 일간신문에 공고하여야 한다. 다만, 토지등소유자 1인이 시행하는 재개발사업의 경우에는 그러하지 아니하다. \par
1. \textasciitilde{} 2. <생략> \par
3. 분양신청기간 \par
4. <생략> \par
③ 대지 또는 건축물에 대한 분양을 받으려는 토지등소유자는 제2항에 따른 분양신청기간에 대통령령으로 정하는 방법 및 절차에 따라 사업시행자에게 대지 또는 건축물에 대한 분양신청을 하여야 한다. \par
제73조(분양신청을 하지 아니한 자 등에 대한 조치) ① 사업시행자는 관리처분계획이 인가ㆍ고시된 다음 날부터 90일 이내에 다음 각 호에서 정하는 자와 토지, 건축물 또는 그 밖의 권리의 손실보상에 관한 협의를 하여야 한다. 다만, 사업시행자는 분양신청기간 종료일의 다음 날부터 협의를 시작할 수 있다. \par
1. 분양신청을 하지 아니한 자 \par
2. \textasciitilde{} 4. <생략> \par
② 사업시행자는 제1항에 따른 협의가 성립되지 아니하면 그 기간의 만료일 다음 날부터 60일 이내에 수용재결을 신청하거나 매도청구소송을 제기하여야 한다. \par
<공익사업을 위한 토지 등의 취득 및 보상에 관한 법률 시행규칙> \par
제54조(주거이전비의 보상) ① 공익사업시행지구에 편입되는 주거용 건축물의 소유자에 대하여는 해당 건축물에 대한 보상을 하는 때에 가구원수에 따라 2개월분의 주거이전비를 보상하여야 한다. <단서 생략> \par
② 공익사업의 시행으로 인하여 이주하게 되는 주거용 건축물의 세입자(무상으로 사용하는 거주자를 포함하되, 법 제78조제1항에 따른 이주대책대상자인 세입자는 제외한다)로서 사업인정고시일등 당시 또는 공익사업을 위한 관계 법령에 따른 고시 등이 있은 당시 해당 공익사업시행지구안에서 3개월 이상 거주한 자에 대해서는 가구원수에 따라 4개월분의 주거이전비를 보상해야 한다. <단서 생략> \par
\end{framedtext}
\end{longquestion}
